\section{Résultats}

\subsection{Oscillations libres}

Avant de mesurer une oscillation, nous avons mesure le rapport de force qu'il faut exercer sur le balancier afin de le placer a un certain angle.
\\
\noindent \begin{minipage}{0.5\textwidth}
\vspace{0.4cm}
\includegraphics[width=\textwidth]{raideur.png}
        \captionof{figure}{Amortissement en fonction du courant}
\label{fig:nature}
\end{minipage}
\hspace{0.05\textwidth}
\begin{minipage}{0.4\textwidth}
        Nous obtenons par consequent le graphique \ref{fig:nature} ci-joint.
        Grace a sa pente d'unite ... nous avons pu directement extraire la valeur C, soit la constante de raideur du ressort utilise grace a la relation \ref{eq:nature} Soit :
\\      C = 123.123
\noindent \end{minipage}
\\
\\
\\
En un second temps, apres avoir montee l'installation pour mesures les oscillations du balencier-spiral (voir \ref{fig:samplesetup}), nous avons donne une position hors-equilibre au balencier et mesurer en fonction du frein imposee a les differences oscillations observable en annexe (voir \ref{fig:oscillations}).

\begin{minipage}[t]{0.5\textwidth}
    \centering
    \includegraphics[width=\linewidth]{lambda.png}
        \captionof{figure}{lambdas en fonction des courants}
\end{minipage}
\hfill
\begin{minipage}[t]{0.46\textwidth}
    \centering
    \includegraphics[width=\linewidth]{lambda_quadratique.png}
	\captionof{figure}{Lambdas en fonction des courants au carre}
\end{minipage}
\\
\\
skhvrb skruv slirvb.
\newpage
\subsection{Expérience 2}

En tracant la phase en fonction de la frequence (voir \ref{fig:phase}), nous pouvons retrouver les frequences propres $f_0$ pour les alimentations a 0.5 en bleu et 0.9A en rouge.
        Les frequences propres resultent du point d'inflexions de leurs courbes correspondantes.
\\
\noindent \begin{minipage}{0.5\textwidth}
\vspace{0.4cm}
\includegraphics[width=\textwidth]{force_phase.png}
	\captionof{figure}{Phase en fonction de la frequence}
\label{fig:phase}
\end{minipage}
\hspace{0.05\textwidth}
\begin{minipage}{0.4\textwidth}
	Ci-joint (\ref{fig:phas}), les frequences propres ont pu s'obtenir a partir de la relation pulsation-frequence decrite dans la section theorie (voir \iref{eq:pulsation-frequence}). Cela nous donne les valeurs a 0.5A de : $f_0$ = 0.518 $+/-$ 0.0000 [Hz]. A 0.9A, nous obtenons en revanche : $f_0$ = 0.528 $+/-$ 0.000 [Hz].
\noindent \end{minipage}

\noindent \begin{minipage}{0.5\textwidth}
\vspace{0.4cm}
\includegraphics[width=\textwidth]{force_amp.png}
	\captionof{figure}{Amplitude(f)}
	\label{fig:amplitude}
\label{fig:nature}
\end{minipage}
\hspace{0.05\textwidth}
\begin{minipage}{0.4\textwidth}
	Comme pour la phase en fonction de la frquence, l'amplitude du balencier en fonction de la frequence nous permet d'obtenir les differents parametres de notre systemes ($\lambda , \omega_0 , \f_0)$. 
	Toutefois, c'est sheitan.
\noindent \end{minipage}
\\
\\
Nous pourrions etablir une moyenne des frequences ou les valeurs des pulsations propres et des pseudo pulsations a partir d'une moyenne. Toutefois, ces donnees sont extraites des fits dont la qualite differe entre celui de la phase (\ref{fig:phase}) et celui de l'amplitude (\ref{fig:amplitude}). Ainsi, il est plus judicieux d'etablir avec plus de certitude qu'a partir de la phase en fonction de la frequence que $\lambda = $. $\omega_0 = $ . Et la par moyenne des deux freqences : $\f_0 = 0.523 +/- 0.00 [Hz]$. 
\\
\\
Enfin, si nous etablissons le rapport $P$ / $P_0$ en fonction de la frequence a differents courants induits dans les freins, nous pouvons observer une courbe plus aplatie (voir \ref{fig:power9}) qu'en le graphique a 0.5 A (voir \ref{fig:power5}). 

\begin{minipage}[t]{0.46\textwidth}
    \centering
    \includegraphics[width=\linewidth]{power_0.5A.png}
	\captionof{figure}{$P$/$P_0$(f) 0.5A}
	\label{fig:power5}
\end{minipage}
\hfill
\begin{minipage}[t]{0.46\textwidth}
    \centering
    \includegraphics[width=\linewidth]{power_0.9A.png}
	\captionof{figure}{$P$/$P_0$(f) a 0.9A}
	\label{fig:power9}
\end{minipage}


